\subsection*{范围：\texttt{\#name}}

\begin{lstlisting}[language=croqtile,caption={范围：\#name}]
foreach i in [128 / #tile]
\end{lstlisting}

前缀 \co@#tile@ 返回 \co@tile@ 的范围。声明为 \co@parallel tile by 8@ 时，\co@#tile@ 就是 \co@8@。

图~\ref{fig:f04-fig_extent} 展示了 \texttt{\#} 前缀的两种用法。

\begin{figure}[htbp]\centering
\includegraphics[width=\linewidth,height=0.56\textheight,keepaspectratio]{f04-fig_extent_light.png}
\caption{前缀 \#tile 返回轴的范围大小，此处为 8；中缀 tile \# i 将分块坐标与块内偏移组合。}
\label{fig:f04-fig_extent}
\end{figure}
